\documentclass[12pt]{article}
\usepackage[utf8]{inputenc}
\usepackage[T1]{fontenc}
\usepackage[margin=1in]{geometry}
\usepackage{setspace}
\usepackage{xcolor}
\usepackage[colorlinks=true, urlcolor=blue!60!black, linkcolor=blue!60!black]{hyperref}
\setlength{\parindent}{2em}
\setlength{\parskip}{1.25em}
\renewcommand{\baselinestretch}{1.0}
\begin{document}
\begin{center}
{\LARGE \textit{Chisholm v. Georgia}, 2 U.S. (2 Dall.) 419 (1793)\par}
\vspace{0.8em}
{\large Case Brief\par}
\end{center}
\vspace{1em}

\section{Memory Jogger}
In the Court's first major constitutional decision, a South Carolina executor sued Georgia on a Revolutionary War debt; the Justices held 4--1, in seriatim opinions, that Article III lets a citizen of one State sue another State without its consent---a ruling the Eleventh Amendment promptly overturned (\href{https://supreme.justia.com/cases/federal/us/2/419/}{Justia}).

\section{Detailed Case Facts}
In 1777, two commissioners authorized to act for the State of Georgia, Thomas Stone and Edward Davies, purchased Revolutionary War supplies for American troops near Savannah from Robert Farquhar, a South Carolina merchant (\href{https://www.fjc.gov/history/cases/cases-that-shaped-the-federal-courts/chisholm-v-georgia}{Federal Judicial Center}). Although Georgia's treasury provided funds to the commissioners, they never paid Farquhar, who died in 1784 with the debt unpaid. Alexander Chisholm, a South Carolina citizen and executor of Farquhar's estate, petitioned the Georgia legislature for payment, which denied his claim in 1789 (\href{https://www.fjc.gov/history/cases/cases-that-shaped-the-federal-courts/chisholm-v-georgia}{Federal Judicial Center}).

\section{Procedural History}
Chisholm first sued Georgia in the U.S. Circuit Court for the District of Georgia (\emph{Farquhar's Executor v. Georgia}). Georgia Governor Edward Telfair and the State's Attorney General refused to submit to jurisdiction, asserting sovereign immunity. In October 1791, Justice James Iredell (riding circuit) and District Judge Nathaniel Pendleton dismissed the case for lack of jurisdiction under the Judiciary Act of 1789 (\href{https://www.fjc.gov/history/cases/cases-that-shaped-the-federal-courts/chisholm-v-georgia}{Federal Judicial Center}).

In 1792, Chisholm refiled an action of \emph{assumpsit} directly in the Supreme Court of the United States under Article III, Section 2, and Section 13 of the Judiciary Act of 1789 (\href{https://supreme.justia.com/cases/federal/us/2/419/}{2 U.S. at 430--431}). Summons was served on Georgia's Governor and Attorney General. When Georgia refused to appear in August 1792, Attorney General Edmund Randolph moved for a conditional default judgment order and writ of inquiry unless Georgia appeared or showed cause at the next Term (\href{https://supreme.justia.com/cases/federal/us/2/419/}{2 U.S. at 429}). Georgia again declined to appear in February 1793, and the Court heard argument solely from Randolph.

\section{Judicial Votes}
The Supreme Court, exercising original jurisdiction, ruled for the plaintiff by a \textbf{4--1} vote and ordered that default judgment be entered unless Georgia appeared or showed cause. The case was argued on February 5, 1793, and decided on February 19, 1793 (\href{https://supreme.justia.com/cases/federal/us/2/419/}{Justia}):

\begin{itemize}
    \item \textbf{Majority (Seriatim Opinions):} Chief Justice Jay and Justices Cushing, Wilson, and Blair each delivered a separate opinion; there was no single opinion of the Court.
    \item \textbf{Concurring Opinions:} None apart from the seriatim majority opinions.
    \item \textbf{Dissenting Opinion:} Justice Iredell dissented alone.
\end{itemize}

Five Justices participated, and the seriatim opinions were delivered in reverse order of seniority.

\section{Holding}
\textbf{Questions Presented:}

\begin{enumerate}
  \item \textbf{Constitutional Suability:} Under Article III, Section 2 of the U.S. Constitution---extending federal judicial power to ``controversies ... between a State and citizens of another State'' and granting the Supreme Court original jurisdiction in cases ``in which a State shall be a party''---may a citizen of one State sue another State as a defendant in the Supreme Court without its consent?
  \item \textbf{Statutory Cognizance and Process:} Under Sections 13 and 14 of the Judiciary Act of 1789, may an action of \emph{assumpsit} for the recovery of money lie against a nonconsenting State, and is service upon the Governor and Attorney General legally sufficient?
\end{enumerate}

\textbf{Holding:} \textbf{Yes (4--1).} Article III, Section 2 expressly extends federal judicial power to controversies between a State and citizens of another State regardless of whether the State is plaintiff or defendant, and vests the Supreme Court with original jurisdiction over such suits (\href{https://supreme.justia.com/cases/federal/us/2/419/}{2 U.S. at 450--453, 466, 469, 479}). Because sovereignty resides in the people of the United States rather than the individual States, common-law sovereign immunity does not bar an out-of-state citizen's \emph{assumpsit} action against Georgia. Service on the Governor and Attorney General was valid, and the Court ordered that default judgment be entered unless Georgia appeared or showed cause by the first day of the next Term (\href{https://supreme.justia.com/cases/federal/us/2/419/}{2 U.S. at 479}).

\section{Analysis of Opinions}
\subsection{Seriatim Majority Opinions (Chief Justice Jay and Justices Cushing, Wilson, and Blair)}

Following the early Court's \emph{seriatim} practice, the four majority Justices each delivered separate opinions combining textual, philosophical, and structural arguments:

\begin{itemize}
  \item \textbf{Plain Text of Article III (Blair, Cushing, Jay, Wilson):} Article III, Section 2 extends judicial power to ``controversies ... between a State and citizens of another State'' and grants original jurisdiction where ``a State shall be a party.'' Rejecting the ratification-era view of Alexander Hamilton (\emph{Federalist No.~81}), James Madison, and John Marshall that States could appear only as plaintiffs (\href{https://www.law.cornell.edu/constitution-conan/amendment-11/historical-background-on-eleventh-amendment}{Cornell LII}), Justice Blair and Justice Cushing emphasized that adjacent clauses (``between two or more States'' and ``between a State ... and foreign states'') necessarily make States defendants, and no plaintiff-only exception appears in the text (\href{https://supreme.justia.com/cases/federal/us/2/419/}{2 U.S. at 450--451, 467}). Chief Justice Jay noted that ``party'' in Article III and Section 13 of the Judiciary Act includes both plaintiffs and defendants; Congress distinguished plaintiffs from defendants for ambassadors in the same section but not for States (\href{https://supreme.justia.com/cases/federal/us/2/419/}{2 U.S. at 477}).
  \item \textbf{Popular Sovereignty Over Feudal Immunity (Wilson and Jay):} Framing the core question as ``do the people of the United States form a Nation?'' (\href{https://supreme.justia.com/cases/federal/us/2/419/}{2 U.S. at 453}), Justice Wilson argued that English royal immunity stemmed from feudal hierarchy and Blackstone's command theory of law, whereas under the U.S. Constitution ``there are citizens, but no subjects,'' and States are artificial persons subordinate to the people (\href{https://supreme.justia.com/cases/federal/us/2/419/}{2 U.S. at 455--458}). ``As to the purposes of the Union, therefore, Georgia is NOT a sovereign state'' (\href{https://supreme.justia.com/cases/federal/us/2/419/}{2 U.S. at 457}). Chief Justice Jay likewise traced sovereignty from the British Crown directly to the people as ``joint tenants in the sovereignty,'' noting that Georgia was simultaneously suing a British creditor and his South Carolina partners as a plaintiff in \emph{Georgia v. Brailsford} and that ``[t]hat rule is said to be a bad one which does not work both ways'' (\href{https://supreme.justia.com/cases/federal/us/2/419/}{2 U.S. at 471--473}).
  \item \textbf{Constitutional Structure and Domestic Tranquility (Wilson, Cushing, Jay):} Wilson reasoned that the Contracts Clause (Article I, Section 10) would be hollowed out if a State could breach its own contracts without judicial accountability (\href{https://supreme.justia.com/cases/federal/us/2/419/}{2 U.S. at 465}). Jay and Cushing stressed that a neutral national tribunal was vital to establish justice and prevent interstate commercial disputes from escalating into conflict (\href{https://supreme.justia.com/cases/federal/us/2/419/}{2 U.S. at 467--469, 474--476}). Addressing whether the United States could likewise be sued, Cushing noted textual differences in Article III, while Jay observed that federal courts have no executive power to enforce judgments against the United States itself (\href{https://supreme.justia.com/cases/federal/us/2/419/}{2 U.S. at 469, 478}).
\end{itemize}

\subsection{Distinctions Among the Seriatim Majority Opinions}

\begin{itemize}
  \item \textbf{Justice Blair} (\href{https://supreme.justia.com/cases/federal/us/2/419/}{2 U.S. at 450--453}) relied strictly on constitutional text, declining to invoke European history or political philosophy and rejecting enforcement objections as premature.
  \item \textbf{Justice Cushing} (\href{https://supreme.justia.com/cases/federal/us/2/419/}{2 U.S. at 466--469}) emphasized textual reciprocity and presciently noted that ``[i]f the Constitution is found inconvenient in practice in this or any other particular, it is well that a regular mode is pointed out for amendment.''
  \item \textbf{Justice Wilson} (\href{https://supreme.justia.com/cases/federal/us/2/419/}{2 U.S. at 453--466}) and \textbf{Chief Justice Jay} (\href{https://supreme.justia.com/cases/federal/us/2/419/}{2 U.S. at 469--479}) advanced sweeping nationalist theories of popular sovereignty.
\end{itemize}

\subsection{Dissenting Opinion (Justice Iredell)}

Justice Iredell dissented alone (\href{https://supreme.justia.com/cases/federal/us/2/419/}{2 U.S. at 429--450}), framing the narrow question as whether an action of \emph{assumpsit} could lie against a State:
\begin{enumerate}
  \item \textbf{Legislative Authority Required:} Federal courts cannot fashion new remedies on their own authority; under the Necessary and Proper Clause, Congress must prescribe the methods of proceeding (\href{https://supreme.justia.com/cases/federal/us/2/419/}{2 U.S. at 432--433, 448--449}).
  \item \textbf{Common-Law Limits Under Section 14:} Because Congress enacted no special remedy against States, Section 14 of the Judiciary Act of 1789 limited federal writs to those ``agreeable to the principles and usages of law'' (\href{https://supreme.justia.com/cases/federal/us/2/419/}{2 U.S. at 433--434}). Reviewing English common law (including \emph{The Bankers' Case}, Lord Somers, and Blackstone), Iredell concluded that no compulsory \emph{assumpsit} action ever lay against the Crown; relief was solely by petition of right as ``a matter of grace, though not upon compulsion'' (\href{https://supreme.justia.com/cases/federal/us/2/419/}{2 U.S. at 437--445}).
  \item \textbf{Dual Sovereignty and Rejection of the Corporation Analogy:} Each State remains ``completely sovereign'' as to all reserved powers (\href{https://supreme.justia.com/cases/federal/us/2/419/}{2 U.S. at 435}). Unlike subordinate municipal corporations created and revocable by a legislature, States predated the Union and derive their authority from their people (\href{https://supreme.justia.com/cases/federal/us/2/419/}{2 U.S. at 447--448}).
  \item \textbf{Constitutional Intimation:} While resting formally on the Judiciary Act, Iredell added an ``extrajudicial'' intimation strongly opposing any constitutional construction permitting ``a compulsive suit against a State for the recovery of money'' (\href{https://supreme.justia.com/cases/federal/us/2/419/}{2 U.S. at 449--450}).
\end{enumerate}

\subsection{Subsequent Developments}

\begin{itemize}

\item \textbf{The Eleventh Amendment:} Congress proposed the \href{https://www.law.cornell.edu/constitution-conan/amendment-11/historical-background-on-eleventh-amendment}{Eleventh Amendment} on March 4, 1794, and States ratified it on February 7, 1795: ``The Judicial power of the United States shall not be construed to extend to any suit in law or equity, commenced or prosecuted against one of the United States by Citizens of another State, or by Citizens or Subjects of any Foreign State.'' Although the Court entered default judgment for Chisholm in 1794, no writ of inquiry was executed; in \href{https://supreme.justia.com/cases/federal/us/3/378/}{\emph{Hollingsworth v. Virginia}, 3 U.S. (3 Dall.) 378 (1798)}, the Court unanimously dismissed all pending suits against States.

\item \textbf{In-State and State-Court Immunity:} \href{https://supreme.justia.com/cases/federal/us/134/1/}{\emph{Hans v. Louisiana}, 134 U.S. 1 (1890)}, barred suits by a State's own citizens in federal court; \href{https://supreme.justia.com/cases/federal/us/527/706/}{\emph{Alden v. Maine}, 527 U.S. 706 (1999)}, barred Article~I private damages suits in a State's own courts; and \href{https://supreme.justia.com/cases/federal/us/587/17-1299/}{\emph{Franchise Tax Board of California v. Hyatt}, 587 U.S. \_\_\_ (2019)}, recognized interstate sovereign immunity in sister-state courts.
  \item \textbf{Officer Suits and Abrogation:} \href{https://supreme.justia.com/cases/federal/us/209/123/}{\emph{Ex parte Young}, 209 U.S. 123 (1908)}, permits prospective injunctive relief against state officers. Congress generally cannot abrogate immunity under Article~I (\href{https://supreme.justia.com/cases/federal/us/517/44/}{\emph{Seminole Tribe}, 517 U.S. 44 (1996)}), but may do so under Section~5 of the Fourteenth Amendment (\href{https://supreme.justia.com/cases/federal/us/427/445/}{\emph{Fitzpatrick v. Bitzer}, 427 U.S. 445 (1976)}) if congruent and proportional (\href{https://supreme.justia.com/cases/federal/us/589/18-877/}{\emph{Allen v. Cooper}, 589 U.S. \_\_\_ (2020)}).

\item \textbf{Arm-of-the-State Test:} On March 4, 2026, in \href{https://supreme.justia.com/cases/federal/us/607/24-1021/}{\emph{Galette v. New Jersey Transit Corp.}, 607 U.S. \_\_\_ (2026)}, a unanimous Court held that NJ Transit is a legally separate corporation responsible for its own judgments---not an ``arm of the State''---and therefore lacks interstate sovereign immunity.

\end{itemize}

\section{Examples: Future Applications}
\begin{enumerate}

\item \textbf{Same-Side (No Immunity Bar): The State as Plaintiff.} A Texas supplier sells generators to Louisiana and then refuses to honor its warranty. Flip the posture of a typical immunity case: if \emph{Louisiana} sues the Texas supplier in federal court, immunity does not apply---creating the one-way asymmetry Chief Justice Jay criticized (\href{https://supreme.justia.com/cases/federal/us/2/419/}{2 U.S. at 473}). The Eleventh Amendment reaches only suits ``commenced or prosecuted against'' a State, so Article III permits federal jurisdiction when the State sues. The Supreme Court has original, nonexclusive jurisdiction over such an action, and a lower federal court would need a separate statutory grant. Georgia likewise sued Samuel Brailsford (a British subject) and his South Carolina partners in the Supreme Court in \emph{Georgia v.\ Brailsford} while claiming immunity in \emph{Chisholm}.

\item \textbf{Same-Side (Relief Available): Prospective Relief Against a State Officer.} After Louisiana's legislature passes a statute cancelling its outstanding contracts with out-of-state suppliers, a Texas supplier sues the state treasurer in federal court, seeking an injunction against enforcement of the statute as a violation of the Contracts Clause. The suit may proceed: under \href{https://supreme.justia.com/cases/federal/us/209/123/}{\emph{Ex parte Young}, 209 U.S.\ 123 (1908)}, a federal court may order a state officer to stop violating federal law prospectively, and the officer cannot shelter behind the State's immunity. This preserves part of the \emph{Chisholm} majority's concern that a State not escape its contracts without judicial accountability. The supplier cannot, however, obtain an order requiring payment of the past debt from the state treasury.

\item \textbf{Opposite-Side (Suit Barred): An Unpaid Contract Debt in Federal Court.} A Texas supplier sues the State of Louisiana in federal court for \$15 million on an unpaid generator contract, alleging breach of contract and a Contracts Clause violation. Under \emph{Chisholm}, the suit proceeds (\href{https://supreme.justia.com/cases/federal/us/2/419/}{2 U.S. at 456}); today, it is barred by the \href{https://www.law.cornell.edu/constitution-conan/amendment-11/historical-background-on-eleventh-amendment}{Eleventh Amendment} and \href{https://supreme.justia.com/cases/federal/us/134/1/}{\emph{Hans}}, and \href{https://supreme.justia.com/cases/federal/us/209/123/}{\emph{Ex parte Young}} does not permit retroactive treasury damages.

\item \textbf{Opposite-Side (Suit Barred): Federal Wage Claims in the State's Own Courts.} State probation officers in Maine sue the State in a Maine state court, seeking back overtime pay owed under the federal Fair Labor Standards Act. They argue that the Eleventh Amendment, which speaks only to ``the Judicial power of the United States,'' says nothing about state courts. Under \href{https://supreme.justia.com/cases/federal/us/527/706/}{\emph{Alden v.\ Maine}, 527 U.S.\ 706 (1999)}, the suit is barred: Congress may not use its Article I powers to subject a nonconsenting State to private damages suits in its own courts. The result rests on the structural view of state sovereignty traced to Justice Iredell's dissent, not on the Amendment's text that the \emph{Chisholm} majority's textualism would have emphasized.

\item \textbf{Fence-Sitter (Close Call): Is a State-Created Hospital Authority an ``Arm of the State''?} A State creates a public hospital authority as a separate corporation empowered to contract and to sue or be sued in its own name. Its charter makes it legally responsible for its own debts, but a later statute requires the State treasury to satisfy any final judgment against it above a specified threshold. An out-of-state vendor sues in federal court for unpaid invoices and seeks a judgment above that threshold. Under \textit{Galette v.\ New Jersey Transit Corp.}\ (2026), separate incorporation presumptively weighs against immunity, while a State's formal legal responsibility for judgments supports arm-of-the-State status. The question is whether the State's mandatory but limited assumption of liability, considered with the authority's independent legal status, rebuts that presumption.

\end{enumerate}

\section{Critique}
\textbf{Textualism Inverted:} The \emph{Chisholm} majority applied literal textualism to Article~III, whereas modern sovereign immunity relies on unwritten structural postulates drawn from Iredell's dissent (which itself rested formally on the Judiciary Act of 1789, \href{https://supreme.justia.com/cases/federal/us/2/419/}{2 U.S. at 449--450}).

\textbf{Diversity vs.\ Immunity Theories:} Critics argue the Eleventh Amendment merely repealed Citizen--State diversity jurisdiction, leaving federal-question suits intact.

\section{Key Quotations}
\begin{quote}
``To the Constitution of the United States, the term SOVEREIGN, is totally unknown. There is but one place where it could have been used with propriety. But even in that place, it would not, perhaps, have comported with the delicacy of those who ordained and established that Constitution. They might have announced themselves `SOVEREIGN' people of the United States. But serenely conscious of the fact, they avoided the ostentatious declaration.''
\end{quote}
\noindent --- Justice James Wilson, \href{https://supreme.justia.com/cases/federal/us/2/419/}{2 U.S. at 454}.

\begin{quote}
``A state, like a merchant, makes a contract. A dishonest state, like a dishonest merchant, wilfully refuses to discharge it. The latter is amenable to a court of justice. Upon general principles of right, shall the former, when summoned to answer the fair demands of its creditor, be permitted, Proteus-like, to assume a new appearance, and to insult him and justice by declaring `I am a Sovereign state?' Surely not.''
\end{quote}
\noindent --- Justice James Wilson, \href{https://supreme.justia.com/cases/federal/us/2/419/}{2 U.S. at 456}.

\begin{quote}
``No such ideas obtain here; at the Revolution, the sovereignty devolved on the people, and they are truly the sovereigns of the country, but they are sovereigns without subjects (unless the African slaves among us may be so called), and have none to govern but themselves; the citizens of America are equal as fellow citizens, and as joint tenants in the sovereignty.''
\end{quote}
\noindent --- Chief Justice John Jay, \href{https://supreme.justia.com/cases/federal/us/2/419/}{2 U.S. at 471--472}.

\begin{quote}
``Every State in the Union, in every instance where its sovereignty has not been delegated to the United States, I consider to be as completely sovereign as the United States are in respect to the powers surrendered. The United States are sovereign as to all the powers of Government actually surrendered: each State in the Union is sovereign as to all the powers reserved.''
\end{quote}
\noindent --- Justice James Iredell, dissenting, \href{https://supreme.justia.com/cases/federal/us/2/419/}{2 U.S. at 435}.

\end{document}
